\section{Discussion}
\subsection{Principal Findings and Interpretation of the Five-Mode Design}
The five-mode evaluation yielded three main observations. Mode C had the shortest completion time among the tested modes. In the prespecified C--E comparison, the complete parameter bundle reduced mean completion time by 1.79 s (8.5\%), with a matched-block bootstrap 95\% CI of 1.10--2.51 s; the direction favored mode C for all three operators and all six objects. Success rate and Raw NASA-TLX descriptively favored mode C, whereas the trajectory-length interval included zero. Taken together, these findings support a within-sample system-level advantage of the complete joint configuration under the tested conditions, rather than an isolated effect of any one parameter channel.

The five modes clarify the scope of this interpretation. Mode B included object identification and manual strategy selection after timing began, so its contrast with mode C represents an automated-workflow comparison rather than an isolated controller comparison. Modes C--E retained the same visual-semantic display throughout the trial. In particular, mode D displayed the same class, mapped strategy, confidence, and bounding box as mode C while retaining fixed parameters; the C--D pattern is therefore consistent with an additional task-level benefit beyond displaying the semantic cue alone. The primary C--E contrast is more specific: both modes used the same class-to-strategy mapping, visual-semantic display, and slave-side impedance update, while mode C additionally updated the master-side haptic gain, force dead zone, gripper speed, and requested grasp-force setting.

Accordingly, the C--E result is consistent with an additional system-level benefit from configuring the haptic-interface and gripper settings together with slave-side impedance. It does not identify the contribution of either added channel or test an interaction between them. The completion-time difference was not accompanied by clear evidence of a shorter master-side geometric path, and the lower mean pause count in mode C remained exploratory because the two modes had the same median and no confirmatory pause analysis was prespecified. The available process measures may be compatible with smoother task progression or fewer corrective interruptions, but they do not establish a specific mechanism.

\subsection{Asynchronous Perception--Control Integration}
A one-shot configuration event does not require semantic perception to operate at the supervisory-loop rate. In the present implementation, 15-fps image acquisition and independent inference exchanged data with the nominal 200-Hz supervisory update through bounded queues. The single-slot frame queue limited stale-image accumulation, and non-blocking polling prevented the supervisory loop from waiting for visual output. This architecture is suited to task-initial configuration rather than continuous visual servoing.

The implementation therefore demonstrates asynchronous, non-blocking integration of low-rate semantic perception into the software architecture. It does not establish hard real-time performance, a measured end-to-end trigger latency, or completion of configuration before contact in every trial.

\subsection{Variation Across Operators and Objects}
The mean paired C--E difference was positive for all three operators and remained positive when each operator was excluded in turn. Across the six objects, mode C was faster than mode E, with descriptive mean reductions from 3.3\% to 13.2\%. These directional checks indicate that the observed aggregate difference was not visibly driven by a single operator or object in the present dataset.

The magnitude varied across objects: it was smaller for the bottle and banana and larger for the paper cup and scissors. Such variation may reflect differences in grasp pose, deformation risk, and orientation-stability requirements. However, each object contributed only four or five matched blocks; the object-level results remain descriptive and cannot establish an object-specific mechanism or generalization to untested objects.

\subsection{Relation to Previous Work and Mechatronic Design Implications}
Previous tele-impedance and variable-impedance approaches commonly adapt remote impedance according to human state, contact force, demonstrations, or task phase \cite{ref06,ref07,ref08,ref09,ref18,ref19,ref20,ref21}. Visual methods have translated object attributes, geometry, and vision--language inputs into stiffness or impedance matrices \cite{ref10,ref11,ref13,ref14}, while automatic switching, shared control, and multi-stiffness interfaces have addressed mode selection and control-authority allocation \cite{ref12,ref15,ref16,ref23,ref24}. Teleoperated grasping has also used tactile deformation and grip-safety information to regulate gripper impedance \cite{ref30}.

The present work complements these approaches with a constrained systems architecture in which one semantic strategy event assigns a consistent operating state across slave mechanical response, operator-facing haptic rendering, and gripper execution. The contribution is thus a cross-subsystem, task-level configuration architecture and its system-level evaluation, rather than a new impedance-control law or a dynamically derived coupling law. Bounded queues, non-blocking polling, one-shot triggering, and strategy locking allow low-rate perception to remain outside the nominal high-rate supervisory path while avoiding repeated intra-trial reconfiguration.

This design principle may be useful in semi-structured applications such as flexible sorting, remote maintenance, and dismantling, where a limited set of recurring object or task classes can be associated with validated operating strategies. Scaling beyond a small class set would require sustained maintenance of the mapping and potentially more than three strategy levels. Additional tactile or contact-state sensing could support post-contact verification or correction without placing the visual process in the high-rate supervisory path.

\subsection{Limitations and Future Work}
The most direct limitation concerns online triggering and timing verification. Formal-trial logs did not retain complete visual-event histories or synchronized strategy-event--contact timestamps. The trigger--contact interval, trigger correctness, fallback use, and completion of configuration before contact therefore could not be reconstructed for individual trials. The 180-image test characterizes recognition only under controlled imaging conditions. The framework should consequently be interpreted as designed for contact-preparatory configuration, not as having verified pre-contact configuration in every trial.

The experimental scope was also limited. The study used repeated measurements from three operators; all were right-handed but operated the left-positioned Omega.7 with the left hand. Only translational master motion was mapped, end-effector orientation was fixed, objects were manually reset within a prescribed region, and the preassigned mode sequences were not completely counterbalanced. The results thus apply to the present device layout, training conditions, motion mapping, placement protocol, object set, and participant sample, and do not support population-level human-factors inference.

The mapping and parameter design impose further attribution limits. The class-to-strategy mapping was manually engineered; a new object requires a reliable detector label and an explicit strategy entry, and the system does not infer unseen physical properties or an appropriate strategy automatically. The three parameter sets were qualitatively screened operating points rather than optimized settings or material-strength limits. Moreover, the C--E comparison changed the haptic-interface and gripper groups together, so channel-specific ablations are required to identify their individual effects. No formal stability or passivity analysis was conducted for the time-varying human-in-the-loop system. The requested Franka Hand \texttt{force} setting was not a measured fingertip force, and the Omega.7 command had no additional software saturation whose status was logged.

Finally, the available outcome and process records limit mechanism analysis. Success and failure were judged in real time by one designated researcher, without systematic video recording or independent re-rating. Failure did not stop the timer immediately, and a dropped object could not be regrasped; this preserved a common terminal endpoint but did not yield stage-specific failure durations. The records also lacked separate mode-B selection times, quantitative contact-quality measures, direct corrective-action measures, and pause-threshold sensitivity analysis. Future studies should add synchronized perception, contact, interface, and selection logs; video-supported outcome assessment; channel-level ablations; direct contact-force and contact-quality measures; active orientation control; less constrained object placement; and a larger, more diverse participant sample.
